\documentclass{article}

\usepackage{amsmath,amssymb}
\usepackage{xurl}
\usepackage[hidelinks]{hyperref}
\setlength{\emergencystretch}{2em}

% Shared commands for notes in docs/paper-gaps/.

\DeclareUnicodeCharacter{2080}{\ifmmode{}_{0}\else\textsubscript{0}\fi}
\DeclareUnicodeCharacter{2081}{\ifmmode{}_{1}\else\textsubscript{1}\fi}
\DeclareUnicodeCharacter{2082}{\ifmmode{}_{2}\else\textsubscript{2}\fi}
\DeclareUnicodeCharacter{2099}{\ifmmode{}_{n}\else\textsubscript{n}\fi}

\newcommand{\Lean}{\textsc{Lean}}
\ExplSyntaxOn
\NewDocumentCommand{\leanid}{m}{
  \ifmmode
    \textsf{\detokenize{#1}}
  \else
    \group_begin:
    \urlstyle{sf}
    \tl_set:Nn \l_tmpa_tl {#1}
    \tl_replace_all:Nnn \l_tmpa_tl {\_} {_}
    \exp_args:NV \path \l_tmpa_tl
    \group_end:
  \fi
}
\ExplSyntaxOff
\providecommand{\tr}{\operatorname{tr}}
\newcommand{\tnleanIssueBase}{https://github.com/LionSR/TNLean/issues/}
\newcommand{\tnleanPrBase}{https://github.com/LionSR/TNLean/pull/}
\newcommand{\tnleanDocBase}{https://sirui-lu.com/TNLean/docs/}
\newcommand{\ghissue}[1]{\href{\tnleanIssueBase#1}{GitHub issue \##1}}
\newcommand{\ghpr}[1]{\href{\tnleanPrBase#1}{PR \##1}}
\newcommand{\doclink}[1]{\href{\tnleanDocBase#1.html}{\path{#1}}}

% Dirac notation and the identity operator, as in blueprint/src/macros/common.tex.
% \providecommand keeps note-local definitions authoritative.
\usepackage{dsfont}
\providecommand{\ket}[1]{|#1\rangle}
\providecommand{\bra}[1]{\langle #1|}
\providecommand{\braket}[2]{\langle #1 | #2 \rangle}
\providecommand{\ketbra}[2]{|#1\rangle\!\langle #2|}
\providecommand{\Id}{\mathds{1}}
\providecommand{\one}{\mathds{1}}
\providecommand{\C}{\mathbb{C}}
\providecommand{\MN}[1]{M_{#1}(\C)}
\providecommand{\MD}{M_D(\C)}

% External referencing for paper sources.  The build helper writes sanitized
% auxiliary files under build/paper-aux/ and excludes duplicated source labels.
\usepackage{xr}

% Import a generated paper auxiliary file with a caller-supplied label prefix.
% The candidate paths support builds from docs/paper-gaps/, the repository root,
% and build/paper-aux/ itself.
\newcommand{\importpaperaux}[2]{%
  \IfFileExists{../../build/paper-aux/#2.aux}{%
    \externaldocument[#1]{../../build/paper-aux/#2}%
  }{%
    \IfFileExists{build/paper-aux/#2.aux}{%
      \externaldocument[#1]{build/paper-aux/#2}%
    }{%
      \IfFileExists{#2.aux}{%
        \externaldocument[#1]{#2}%
      }{}%
    }%
  }%
}

% General external reference macros
\newcommand{\paperref}[2]{\ref{#1-#2}}
\newcommand{\papereqref}[2]{(\ref{#1-#2})}

% CPSV16 source (1606.00608 / MPDO-22-12-17-2.tex) external document and shortcuts
\importpaperaux{cpsv16-}{1606.00608/MPDO-22-12-17-2-tnlean}
\newcommand{\cpsvref}[1]{\paperref{cpsv16}{#1}}
\newcommand{\cpsveqref}[1]{\papereqref{cpsv16}{#1}}

% Verdict marker: \gapnote{<kind>}{<status>}. Kind is the policy
% classification (clarification, local-correction, scope-restriction,
% unfaithful, false-source, open-gap); status is open, wip (elimination
% underway), resolved, or historical. Machine-read by the site index;
% prints a verdict line.
\newcommand{\gapnote}[2]{\par\noindent\textbf{Verdict:} #1 (#2).\par}


\title{Per-block RFP Isometry Form and Cross-Block Orthogonality in CPSV16}
\author{}
\date{2026-06-14}

\begin{document}
\maketitle

\section{Source statement}

The pure-state RFP characterization in Section~3 of
\cite{Cirac2016MPDO_arXiv} writes a canonical-form RFP tensor as
\begin{align}
  A^i
    =
  \bigoplus_{j=1}^g\bigoplus_{q=1}^{r_j}
  \mu_{j,q}X_{j,q}\Lambda_jU_j^iX_{j,q}^{-1},
  \tag{\path{III_CFI_RFP}}
\end{align}
where the matrices $\Lambda_j$ are diagonal, positive, and trace-normalized.
The same theorem requires a joint isometry condition
\begin{align}
  \sum_{i=1}^s
  (U_j^i)_{\alpha,\beta}\,
  \overline{(U_{j'}^i)_{\alpha',\beta'}}
  =
  \delta_{j,j'}\delta_{\alpha,\alpha'}\delta_{\beta,\beta'}.
  \tag{\path{eq:III_isometry}}
\end{align}
The factor $\delta_{j,j'}$ is a cross-block orthogonality condition: for
$j\ne j'$, the physical-index isometries associated with the two normal-tensor
blocks are mutually orthogonal.  Corollary~III.cor3 invokes this same isometry
condition for the BNT elements.\footnote{Local source:
\path{Papers/1606.00608/MPDO-22-12-17-2.tex:543--554} and
\path{Papers/1606.00608/MPDO-22-12-17-2.tex:583--589}.}

\section{Formal boundary}

The current formal theorem applies the single-block isometry form separately to
each block of an indexed family.\footnote{Formal declaration:
\leanid{MPSTensor.rfp_nt_structural_full_blocks} in
\doclink{TNLean/MPS/RFP/StructuralFull}.  Introduced in \ghpr{2785}.}  For each
block $k$ it gives
\begin{align}
  A_k^i
    &=
  X_k\Lambda_kU_k^iX_k^{-1},\\
  \sum_i (U_k^i)^\dagger U_k^i
    &= I.
\end{align}
There is no conclusion relating $U_k$ and $U_\ell$ for $k\ne \ell$.  Thus the
formal theorem records a separate isometry condition for each block, but it does
not record the cross-block equations
\begin{align}
  \sum_i
  (U_k^i)_{\alpha,\beta}\,
  \overline{(U_\ell^i)_{\alpha',\beta'}}
  =
  0
  \qquad (k\ne \ell),
\end{align}
which are part of the source condition \path{eq:III_isometry}.

\section{Elimination plan}

A source-faithful version of Corollary~III.cor3 should be stated for a
BNT-family or whole-tensor canonical-form datum, not merely for independent
blocks.  Its conclusion should include a joint family $U$ satisfying
\begin{align}
  \sum_i
  (U_k^i)_{\alpha,\beta}\,
  \overline{(U_\ell^i)_{\alpha',\beta'}}
  =
  \delta_{k,\ell}\delta_{\alpha,\alpha'}\delta_{\beta,\beta'}.
\end{align}
Such a statement must account for the common physical space into which all
blocks embed; in particular, the joint isometry condition has the dimensional
content that the block-pair spaces fit orthogonally inside the physical space.

\section{Conclusion}

The current theorem is a scope restriction.  It is a useful per-block
consequence of the single-block RFP isometry form, but it is not the full
isometry assertion of Corollary~III.cor3 because it omits the cross-block
$\delta_{j,j'}$ orthogonality equations.

\bibliographystyle{alpha}
\bibliography{references}

\end{document}
